\section*{Anexo J. Consideraciones éticas}
\phantomsection
\addcontentsline{toc}{subsection}{Anexo J. Consideraciones éticas}
\label{anexo:etica}
\setcounter{table}{0}
\renewcommand{\thetable}{J.\arabic{table}}
\renewcommand{\theHtable}{J.\arabic{table}}
% Asesoria del 20.11.2026 y anuncio de Canvas de la Direccion de Computacion (lo paso Fabian Alvarado el 20.11.2026):
% "Anexo de Consideraciones Eticas" en la entrega final de PFC2, que mide dos competencias. COM4_4.2: emitir juicios
% informados, argumentados y justos con base en principios eticos, el bien comun y el desarrollo sostenible (emite un
% juicio logico orientado al bien comun a partir del analisis de una situacion etica problematica; aplica principios
% eticos -responsabilidad, justicia, sostenibilidad, seguridad- en el analisis). CS4.4.1: reconocer sus
% responsabilidades profesionales con base en principios legales y eticos (las identifica segun normas -proteccion de
% datos, propiedad intelectual, privacidad-; evalua el impacto de sus decisiones sobre las personas, las organizaciones
% y la sociedad). "En caso no aplique, pongan el comentario respectivo."

Este anexo desarrolla, para este proyecto, las dos competencias que pide
la Dirección de Computación: COM4\_4.2, emitir juicios éticos
informados y orientados al bien común, y CS4.4.1, reconocer las
responsabilidades profesionales con base en principios legales y
éticos. En resumen: el proyecto no trató datos de personas, y su
cuestión ética principal está en el uso de sus resultados, porque un
control aprendido en simulación no debe llegar a un semáforo real como
si su mejora estuviera probada en la calle.

\subsection*{Un juicio ético sobre el uso de los resultados (COM4\_4.2)}

\noindent\textbf{La situación.} Con los resultados del
Capítulo~\ref{cap:resultados}, una municipalidad podría proponer
instalar el modelo en el Óvalo Vasco: en la simulación, baja el retraso
un 23 a 24\,\% frente al plan actual, el que opera hoy
(Tabla~\ref{tab:lima-resultado}). La decisión afectaría a miles de
personas por hora que no participan en ella, y la propia tesis da cuatro
razones de cautela:
\begin{itemize}
  \item \textbf{La mayor parte de la mejora no viene del aprendizaje.}
    En la misma simulación, con la demanda publicada, el fijo ajustado
    al segundo, un plan fijo sin detectores, baja el retraso un 15 a
    16\,\% frente al plan actual, y el actuado ajustado iguala al modelo.
    Con la demanda alta, el fijo ajustado al segundo supera al modelo.
  \item \textbf{La mejora es de una simulación.} La demanda es en parte
    un supuesto, no se simulan peatones y ningún caso se comprobó con
    tiempos de viaje medidos.
  \item \textbf{El modelo necesita ver todo el acceso,} cerca de 1~km en
    el nudo. En la calle eso pide un detector de video o de radar
    (Sección~\ref{sec:costo-escala}), y el video capta placas y rostros.
  \item \textbf{La recompensa cuenta vehículos, no personas.} Un bus con
    30 pasajeros pesa en ella lo mismo que un auto con uno.
\end{itemize}
% fuente: secciones/tablas/ovalo_resultado_fijo.tex (Tabla tab:lima-resultado): demanda publicada, plan actual 99.5 y
% 100.6 s/veh, fijo ajustado al segundo 84.1 y 84.5 (-15.5 y -16.0 %), actuado ajustado 76.4 y 76.5, modelo de orden
% fijo 76.4 y 76.6 (frente al plan actual, -23.2 y -23.9 %); demanda alta, fijo 133.5 y 127.5 frente al modelo 136.8 y
% 131.6 (+2.5 y +3.2 %); limites: conclusiones.tex ("No se puede afirmar que los casos de Lima predigan el retraso
% real"); acceso observado: capitulo5.tex, sec:costo-escala ("unos 590 m en las entradas del corredor y 1 km en el
% nudo"; detectores de video y radar en la estimacion de costo); recompensa: capitulo3.tex, eq:recompensa-u4 (suma
% sobre vehiculos)

\noindent\textbf{El juicio.} No se debe recomendar instalar el modelo tal
como sale de esta tesis. Con la evidencia disponible, lo que sirve al bien
común es reajustar primero el plan fijo del nudo con conteos actuales:
da la mayor parte de la mejora, no expone a nadie a un control que no se
probó en la calle, no necesita cámaras y es una práctica conocida de la
ingeniería de tránsito. El modelo se justificaría después, como un
piloto de estudio, donde aporte lo que los controles comunes no dan
---en esta tesis, el corredor de Benavides, gracias a los mensajes entre
cruces (Sección~\ref{sec:lima-corredor})---, y con cuatro condiciones:
\begin{enumerate}
  \item \textbf{Seguridad.} El control conserva los estados del
    semáforo, el ámbar, el todo en rojo y los verdes mínimos de los
    peatones, y ante cualquier falla vuelve solo al plan fijo.
  \item \textbf{Justicia.} Se mide el retraso de cada entrada y por
    persona, y la espera de los peatones, y se fija de antemano cuánto
    puede empeorar como máximo cada entrada.
  \item \textbf{Privacidad.} Si se usan cámaras, el equipo se queda
    solo con conteos y velocidades, sin guardar imágenes.
  \item \textbf{Honestidad.} El piloto se compara con el plan fijo
    reajustado, no solo con el plan actual, y su resultado se publica
    aunque sea desfavorable.
\end{enumerate}

\noindent\textbf{Los principios en el análisis.} La
Tabla~\ref{tab:etica-principios} muestra qué pide cada principio en esta
situación, qué hace ya la tesis y qué faltaría para llevar el control a
la calle.

\begin{table}[!htbp]
\setstretch{1}
\centering
\footnotesize
\setlength{\tabcolsep}{3pt}
\begin{tabular}{>{\raggedright\arraybackslash}p{2.5cm}>{\raggedright\arraybackslash}p{3.0cm}>{\raggedright\arraybackslash}p{5.0cm}>{\raggedright\arraybackslash}p{3.9cm}}
\toprule
Principio & Qué pide aquí & Qué hace la tesis & Qué falta para la calle \\
\midrule
Responsabilidad y honestidad & No prometer más de lo que la evidencia muestra & Compara con referencias fuertes, el fijo ajustado al segundo y el actuado ajustado, no solo con el plan actual; escribe sus reglas antes de entrenar, evalúa dos veces e informa también los veredictos EMPATA y FALLA (Sección~\ref{sec:diseno-segundo}); declara sus límites & Medir el efecto en la calle, con conteos y tiempos de viaje \\
\specialrule{0.25pt}{1.5pt}{1.5pt}
Seguridad & Que ningún control dé verde a movimientos que chocan ni quite tiempo a los peatones & El modelo no elige qué luces prender: solo decide si mantiene el verde o pasa al estado siguiente, con ámbar y con verdes mínimo y máximo (Sección~\ref{sec:formulacion}). El todo en rojo del nudo y los verdes mínimos de los peatones del corredor son iguales en todos los controles & Una revisión de seguridad, que SUMO no hace (tampoco simula peatones), y la vuelta automática al plan fijo \\
\specialrule{0.25pt}{1.5pt}{1.5pt}
Justicia & Repartir la espera con equidad entre calles y entre personas & Con el orden de los estados fijo y un verde máximo, ninguna entrada se queda sin verde más de un ciclo; el retraso total cuenta también a los que esperan para entrar a la red (Sección~\ref{sec:plan-resultados}) & Medir por entrada y por persona, y a los peatones: la recompensa cuenta vehículos \\
\specialrule{0.25pt}{1.5pt}{1.5pt}
Sostenibilidad & Que el beneficio no traiga más autos ni más gasto & Entrenar toma horas en un computador de escritorio sin GPU, y el modelo se compara con controles que no necesitan detectores & Medir las emisiones, que la tesis no analiza (Tabla~\ref{tab:metas}), y vigilar si la calle atrae más autos \\
\bottomrule
\end{tabular}
\caption[Consideraciones éticas: los principios en el uso de los resultados]{Los principios éticos aplicados a la situación de este anexo: qué pide cada uno, qué hace ya la tesis y qué faltaría para llevar el control a un semáforo real.}
\label{tab:etica-principios}
\end{table}
% fuente: comparacion con referencias fuertes, reglas antes de entrenar y dos evaluaciones: capitulo3.tex,
% sec:diseno-segundo, y capitulo5.tex, sec:caso-lima ("reglas escritas antes de entrenar"); accion de los agentes
% (mantener o pasar a la fase siguiente, ambar de 3 s, verde de 10 a 60 s): capitulo3.tex, sec:formulacion; todo en
% rojo del nudo igual en todos los controles: anexo_d_osm.tex ("Todo en rojo"); verdes minimos de los peatones en el
% corredor, "Vale para todos los controles": capitulo5.tex, sec:lima-corredor; espera para entrar a la red en el
% retraso total: capitulo5.tex, sec:plan-resultados; computo: capitulo5.tex, sec:costo-escala; emisiones: Tabla
% tab:metas ("una meta de emisiones de CO2, que se retiro porque esa metrica no se uso en el analisis")

\subsection*{Responsabilidades profesionales e impacto (CS4.4.1)}

\noindent\textbf{Responsabilidades según normas y principios.} La
Tabla~\ref{tab:etica-normas} reúne las normas que tocan a este proyecto
como trabajo de ciencia de la computación y cómo se cumplen. Donde una
norma no aplica a lo hecho, se dice, y se indica cuándo aplicaría.

\begin{table}[!htbp]
\setstretch{1}
\centering
\footnotesize
\setlength{\tabcolsep}{3pt}
\begin{tabular}{>{\raggedright\arraybackslash}p{2.2cm}>{\raggedright\arraybackslash}p{4.0cm}>{\raggedright\arraybackslash}p{8.2cm}}
\toprule
Ámbito & Norma o principio & En este proyecto \\
\midrule
Protección de datos personales y privacidad & Ley N.\textsuperscript{o}~29733 \citep{Peru2011_Ley29733} y su reglamento \citep{MINJUSDH2024_Reglamento}; con cámaras, la directiva de videovigilancia \citep{ANPD2020_Videovigilancia}; principio 1.6 de la ACM \citep{ACM2018_Code} & \textbf{No aplica a lo hecho:} no se trataron datos de personas. Los vehículos son simulados; la demanda sale de volúmenes publicados; la red, de OpenStreetMap y de Google Maps, y en el sitio se midieron con cronómetro los tiempos del semáforo. \textbf{Aplicaría al instalar el modelo:} si se observa con video, el equipo debe quedarse solo con conteos y velocidades, sin guardar imágenes, y la entidad que lo opere responde por esos datos. \\
\specialrule{0.25pt}{1.5pt}{1.5pt}
Propiedad intelectual & Decreto Legislativo N.\textsuperscript{o}~822, Ley sobre el Derecho de Autor \citep{Peru1996_DL822}; las licencias de los datos y del software; principio 1.5 de la ACM & Se cita cada fuente: métodos, datos y volúmenes publicados. Los datos de OpenStreetMap tienen licencia ODbL 1.0: piden reconocer la fuente, lo que se hace en la Sección~\ref{sec:caso-lima}, y que las redes derivadas de ellos, publicadas en el repositorio del proyecto, se compartan con la misma licencia \citep{OSM_Copyright}. Google permite mostrar sus imágenes en un impreso con la atribución junto a la imagen y sin alterarla \citep{Google_GeoGuidelines}: la llevan las leyendas de las dos figuras satelitales, que solo se anotaron. SUMO (licencia EPL-2.0) y las bibliotecas de Python (licencias libres) se usaron sin modificarlas. Los algoritmos se programaron a partir de sus artículos, que se citan. \\
\specialrule{0.25pt}{1.5pt}{1.5pt}
Ética profesional & Código de Ética y Conducta Profesional de la ACM \citep{ACM2018_Code} & Principios 1.2, evitar el daño; 1.3, ser honesto; 1.4, ser justo; 2.5, evaluar el sistema y sus riesgos de forma completa, y 3.7, cuidar en especial los sistemas que pasan a ser infraestructura de la sociedad, como los semáforos. Se aplican en el juicio anterior y en la Tabla~\ref{tab:etica-impacto}. \\
\specialrule{0.25pt}{1.5pt}{1.5pt}
Normas de tránsito & Manual de Dispositivos de Control de Tránsito Automotor \citep{MTC2024_Manual} & Un control instalado tendría que cumplirlo y contar con la autorización de la autoridad que administra esos semáforos. La tesis no lo verifica: solo compara controles en simulación. \\
\bottomrule
\end{tabular}
\caption[Consideraciones éticas: responsabilidades según normas y principios]{Responsabilidades profesionales del proyecto según normas legales y principios éticos, y cómo se atienden. Donde una norma no aplica a lo hecho, se indica cuándo aplicaría.}
\label{tab:etica-normas}
\end{table}
% fuente: datos usados: capitulo3.tex, sec:validacion-realidad (tiempos del nudo medidos con cronometros; demanda
% publicada, sin aforos), conclusiones.tex (red de OpenStreetMap y medidas en Google Maps); Ley 29733 (2011) y su
% reglamento, D.S. 016-2024-JUS (publicado el 30.11.2024, deroga el D.S. 003-2013-JUS); directiva aprobada por la
% R.D. 02-2020-JUS/DGTAIPD; ODbL: openstreetmap.org/copyright (atribucion y "share-alike" de las bases de datos
% derivadas); repositorio publico: anexos.tex, tab:programas; Google: Geo Guidelines (atribucion a Google y a sus
% proveedores de datos junto a la imagen; anotar si, alterar o quitar la atribucion no); SUMO 1.25.0, EPL-2.0 (con
% GPL-2.0 o posterior como licencia secundaria); NumPy, SciPy, Matplotlib y PyTorch: licencias libres tipo BSD o PSF;
% algoritmos programados sin bibliotecas de aprendizaje por refuerzo: introduccion.tex ("Aporte propio"); manual del
% MTC aprobado por la R.D. 26-2024-MTC/18 (16.10.2024), que deroga la R.D. 16-2016-MTC/14

\noindent\textbf{Impacto de las decisiones.} La
Tabla~\ref{tab:etica-impacto} evalúa el impacto de las decisiones del
proyecto sobre las personas, las organizaciones y la sociedad.

\begin{table}[!htbp]
\setstretch{1}
\centering
\footnotesize
\setlength{\tabcolsep}{3pt}
\begin{tabular}{>{\raggedright\arraybackslash}p{2.4cm}>{\raggedright\arraybackslash}p{2.8cm}>{\raggedright\arraybackslash}p{4.6cm}>{\raggedright\arraybackslash}p{4.6cm}}
\toprule
 & Beneficio posible & Riesgo & Qué decidió el proyecto \\
\midrule
Personas: conductores, pasajeros, peatones y vecinos & Menos espera en la hora punta & Que la mejora media oculte una entrada o a peatones que esperan más; un control que no se probó en la calle & Mantener los estados, el ámbar, el todo en rojo y los verdes de los peatones en todos los controles, y dar el retraso con su intervalo de confianza y sus límites \\
\specialrule{0.25pt}{1.5pt}{1.5pt}
Organizaciones: municipalidades, operadores de semáforos y la universidad & Comparar controles antes de invertir & Pagar detectores por una mejora que un plan fijo reajustado ya da: en el corredor, con la demanda publicada, el tiempo ahorrado frente al fijo vale menos que el detector más barato (Sección~\ref{sec:costo-escala}) & Comparar con el fijo ajustado al segundo y con el actuado ajustado, y estimar el costo y el beneficio \\
\specialrule{0.25pt}{1.5pt}{1.5pt}
Sociedad: la ciudad y el ambiente & Menos congestión y, quizá, menos emisiones & Más autos atraídos por calles más fluidas; un bien público manejado por decisiones automáticas, sin un responsable claro; menos peso para cada pasajero de bus que para quien va en auto & No afirmar lo que no se midió (emisiones, peatones, la calle real) y describir los programas, las reglas y los datos de cada caso (Anexos~D, H e~I) para que otros los revisen \\
\bottomrule
\end{tabular}
\caption[Consideraciones éticas: impacto de las decisiones del proyecto]{Impacto de las decisiones del proyecto sobre las personas, las organizaciones y la sociedad: el beneficio posible, el riesgo y lo que decidió el proyecto.}
\label{tab:etica-impacto}
\end{table}
% fuente: costo y beneficio del corredor: capitulo5.tex, sec:costo-escala ("Frente al fijo, que no necesita
% detectores, el valor ahorrado con la demanda publicada no cubre ni las espiras, el detector mas barato");
% intervalos de confianza: capitulo3.tex, sec:estadistica

\noindent\textbf{Lo que no aplica.} El proyecto no tuvo participantes
humanos: no hubo encuestas, entrevistas ni pruebas con personas, por lo
que no requirió consentimiento informado. Tampoco trató datos
personales (Tabla~\ref{tab:etica-normas}).
